\section{Methodology}
\label{sec:methodology}

We designed experiment h-e1 as a \textbf{MUST\_WORK existence check} --- a protocol that gives projection-only LoRA every possible advantage and treats failure as an unambiguous signal rather than a tuning problem. If LoRA is going to work on Mamba for classification at all, it should work under the conditions we construct.

The key diagnostic insight: we measured \textbf{per-epoch accuracy alongside loss} at every training epoch. Loss curves can look plausible --- oscillating, non-zero, apparently ``doing something'' --- while accuracy reveals complete learning failure. Three consecutive identical accuracy values across epochs cannot be rationalized as slow convergence or unlucky initialization.

\begin{figure}[t]
\centering
% Figure 1: Mamba gradient path diagram
% [Placeholder: selective SSM scan kernel between classification head and LoRA matrices]
\caption{Mamba gradient path diagram. The selective SSM scan kernel sits between the classification head and the LoRA weight matrices in \texttt{in\_proj}/\texttt{out\_proj}/\texttt{x\_proj} across all 24 layers. The dashed arrow from classification head through SSM scan to LoRA matrices indicates the non-discriminative gradient path.}
\label{fig:gradient_path}
\end{figure}

\begin{figure}[t]
\centering
% Figure 2: Two-panel accuracy curve
% [Placeholder: SST-2 flat at 0.5092, MNLI degrading from 0.3463]
\caption{Per-epoch accuracy curves. Left: SST-2 accuracy across epochs (flat line at 0.5092; zero-shot reference at 0.4908; gate criterion at 0.70). Right: MNLI accuracy (0.3463 zero-shot $\to$ 0.3326 epoch~1 $\to$ 0.3234 epoch~2).}
\label{fig:accuracy_curves}
\end{figure}

\subsection{Mamba Architecture: Selective State Space Scan and Gradient Flow}

Mamba-130m (\texttt{state-spaces/mamba-130m-hf}) is a 24-layer causal language model with hidden dimension $d_\text{model} = 768$, inner dimension $d_\text{inner} = 1536$ (expansion factor 2), state dimension $d_\text{state} = 16$, and discrete-time rank $dt\_rank = 48$. Each layer contains a MambaBlock with the following \texttt{nn.Linear} projection layers: \texttt{in\_proj} (768 $\to$ 3072), \texttt{out\_proj} (1536 $\to$ 768), \texttt{x\_proj} (1536 $\to$ 112), and \texttt{dt\_proj} (48 $\to$ 1536).

The MambaBlock also contains a \texttt{conv1d} layer and, critically, the \textbf{selective SSM scan}: the computation $h_t = \bar{A}(\Delta, A) \cdot h_{t-1} + \bar{B}(\Delta, B) \cdot x_t$, $y_t = C \cdot h_t$, implemented as a custom CUDA parallel associative scan (\texttt{selective\_scan\_cuda}).

The gradient path for a classification loss placed on the final hidden state runs: classification head $\to$ final hidden state $\to$ (through scan and projection layers of all 24 layers) $\to$ LoRA weight matrices. The SSM scan sits \emph{between} the classification head and the LoRA targets in every layer. For the model's pretraining task (next-token prediction), this backward path through the scan is well-exercised. For a randomly-initialized classification head computing sequence-level predictions, the backward path is novel --- and as our experiments demonstrate, apparently ineffective at communicating task-discriminative signal.

\subsection{LoRA Application to Mamba-130m}

We applied LoRA using the HuggingFace PEFT library (v0.9+) with the configuration in Table~\ref{tab:lora_config}, matching the community reference \citep{alxndrtl2024mambapeft}.

\begin{table}[h]
\centering
\caption{LoRA configuration for Mamba-130m (Condition A).}
\label{tab:lora_config}
\begin{tabular}{@{}lll@{}}
\toprule
\textbf{Hyperparameter} & \textbf{Value} & \textbf{Rationale} \\
\midrule
Target modules & \texttt{in\_proj}, \texttt{out\_proj}, \texttt{x\_proj} & Community reference configuration \\
Rank ($r$) & 8 & Standard for classification fine-tuning \\
Alpha ($\alpha$) & 16 & $2r$; standard scaling factor \\
Dropout & 0.05 & Standard regularization \\
Trainable parameters & 1,484,288 (1.14\%) & Confirmed by PEFT \texttt{print\_trainable\_parameters()} \\
\bottomrule
\end{tabular}
\end{table}

The classification head is a single \texttt{nn.Linear(768, num\_labels)} layer applied to the last-token hidden state from the final MambaBlock. The head is randomly initialized.

\subsection{GLUE Fine-Tuning Setup}

\begin{table}[h]
\centering
\caption{Training configuration.}
\label{tab:training_config}
\begin{tabular}{@{}ll@{}}
\toprule
\textbf{Hyperparameter} & \textbf{Value} \\
\midrule
Optimizer & AdamW \\
Learning rate & $3 \times 10^{-4}$ \\
Weight decay & 0.01 \\
Batch size & 32 \\
Epochs & 3 \\
LR schedule & Linear warmup (6\%) + linear decay \\
Seed & 42 \\
Training samples & 4,000 (subsampled) \\
Max sequence length & 128 tokens \\
Tokenizer & GPT-NeoX-20B (50k vocab) \\
\bottomrule
\end{tabular}
\end{table}

We fine-tune separately on SST-2 and MNLI. \textbf{MUST\_WORK gate criterion:} SST-2 accuracy $> 0.70$ after three epochs. The conservatism is deliberate: the community reference reports 90--92\%, but a 70\% gate gives LoRA a 20-percentage-point margin for training setup differences.

\textbf{Why 3 epochs, 4,000 samples:} An existence check detects whether a method can learn at all, not its ceiling performance. Three epochs of 4,000 samples is sufficient for LoRA to converge on SST-2 with transformer architectures --- success, if achievable, would be visible.

\subsection{Diagnostic Methodology}

Our core design choice was \textbf{per-epoch accuracy measurement}, not just loss monitoring. Loss curves have multiple patterns consistent with either learning or non-learning. Accuracy is the authoritative criterion: three consecutive identical values across epochs eliminate all alternative explanations (slow convergence, lucky initialization plateau, seed sensitivity, evaluation bug).

We evaluate accuracy on the full validation split at the end of each epoch: SST-2 validation (872 examples), MNLI matched validation (9,815 examples). Before any fine-tuning, we evaluate the base Mamba-130m model on GLUE tasks using \texttt{lm-evaluation-harness} \citep{eval_harness} with multiple-choice prompting to establish the zero-shot reference.

\textbf{Implementation correctness verification:} Before accepting any result, we verify that (1) LoRA keys exist in the model's state dictionary at expected module paths, (2) LoRA weight matrices are non-zero after training, and (3) the trainable parameter count matches the expected value. All three checks pass in our experiments.
